% from ega4-21c-fr.tex
\oldpage[IV-4]{308}
que $j:U\to X$ est affine implique que $U^0$ s'identifie à $U$ et $j^0$ à $j$,
autrement dit que $j^0(U^0)=U$; or, cela contredit la conclusion de (21.12.6) qui,
sous l'hypothèse de dimension faite implique que $j^0$ doit être surjectif.
\end{proof}

\begin{env}[21.12.8]
\label{IV.21.12.8}
Soient $A$ un anneau local noethérien, $Y=\operatorname{Spec}(A)$, $y$ l'unique point
fermé de $Y$, $Y'=Y-\{y\}$. Considérons la condition suivante :
\begin{itemize}
  \item[\textup{(W)}] \emph{Pour tout ouvert $U$ de $Y$ contenu dans $Y'$, ne
  contenant aucune composante irréductible de $Y'$ et tel que l'immersion canonique
  $U\to Y'$ soit affine, $U$ est lui-même un ouvert affine.}
\end{itemize}
Cette condition est entraînée par la suivante :
\begin{itemize}
  \item[\textup{(W')}] \emph{Pour toute partie fermée $T$ de $Y$ dont toute composante
  irréductible est de codimension $1$ dans $Y$, et tel que pour toute composante
  irréductible $Y_i$ de $Y$, on ait $Y_i\cap T\ne\{y\}$, l'ouvert $U=Y-T$ est affine.}
\end{itemize}
En effet, si (W') est vérifiée et si l'ouvert $U$ de $Y$ vérifie l'hypothèse de (W), il
résulte de (21.12.7) qu'aucune composante irréductible de $T=Y-U$ ne peut être de
codimension $\geq2$; comme par ailleurs $U$ ne contient aucune des composantes
irréductibles $Y_i-\{y\}$ de $Y'$, la condition (W') montre que $U$ est affine.

On notera que la condition (W') se simplifie lorsque $Y$ est \emph{irréductible} et est
alors équivalente à la suivante :
\begin{itemize}
  \item[\textup{(W'')}] \emph{Pour toute partie fermée irréductible $T$ de $Y$, de
  codimension $1$ dans $Y$, $Y-T$ est un ouvert affine.}
\end{itemize}
En effet, il est clair que (W') entraîne (W'') lorsque $Y$ est irréductible, et la
réciproque est vraie en considérant les composantes irréductibles de $T$ et en notant
que l'intersection d'un nombre fini d'ouverts affines est un ouvert affine
$(\mathbf I, 5.5.6)$.
\end{env}

\begin{env}[21.12.9]
\label{IV.21.12.9}
\emph{Exemples}. ---
Si $A$ est un anneau local noethérien \emph{factoriel}, il vérifie (W') (et
\emph{a fortiori} (W)), puisque tout idéal premier de hauteur $1$ est principal
$(\mathbf I, 1.3.6)$. Mais il y a des anneaux locaux noethériens non factoriels
vérifiant (W''), par exemple ceux de dimension $\leq1$: en effet, on a remarqué dans la
démonstration de (21.12.6, (i)) que \emph{tous} les ouverts de $Y$ sont alors affines.
On peut d'autre part prouver en utilisant la théorie de la dualité locale (chap.~III,
3\textsuperscript{e} partie) que tout anneau local noethérien de dimension $2$ vérifie
(W').
\end{env}

L'intérêt de la condition (W) réside dans le résultat suivant :

\begin{proposition}[21.12.10]
\label{IV.21.12.10}
Soient $X$, $Y$ deux préschémas localement noethériens, $g:X\to Y$ un morphisme,
$y$ un point fermé de $Y$, $Y'=Y-\{y\}$, $X'=g^{-1}(Y')$; supposons que le morphisme
$g':X'\to Y'$ restriction de $g$ soit une immersion ouverte, et que l'anneau local
$\mathcal O_{Y,y}$ vérifie la condition (W) (21.12.8). Alors, pour toute composante
irréductible $Z$ de $X_y=g^{-1}(y)$, ou bien $Z$ est de codimension $\leq1$ dans $X$,
ou bien son point générique est isolé dans $Z$. Si $Z$ est localement de type fini sur
$k(y)$, la seconde alternative implique que $Z$ est réduit à un seul point.
\end{proposition}

\begin{proof}
La dernière assertion de l'énoncé résulte de ce que $Z$ est alors un préschéma de
Jacobson (10.4.7), et comme l'ensemble des points fermés de $Z$ est alors dense dans
$Z$, le point générique de $Z$ ne peut être isolé dans $Z$ que si $Z$ est réduit à un
seul point.

